\documentclass[12pt, letterpaper]{article}
\usepackage{graphicx}
\usepackage[margin=1in]{geometry}
\usepackage{amsmath, amssymb, amsthm, mathrsfs}
\usepackage{fancyhdr}
\usepackage{tikz}
\usepackage{enumerate}
\usepackage{cancel}
\usetikzlibrary{positioning}
\usepackage[utf8]{inputenc}
\usepackage{xcolor}

\title{HW 1.1}
\author{Your Name}
\date{\today}

\pagestyle{fancy}
\renewcommand{\headrulewidth}{0pt}
\renewcommand{\footrulewidth}{0pt}
\fancyhf{}
\rhead{
    Zack Abdirashid\\
    1/28/26 \\
    MATH 402
}
\rfoot{Page \thepage}

\begin{document}

\paragraph{\underline{2} \\ \\} 

\textbf{1.} \textbf{(a)} \text{Let $p_{1},p_{2},\ldots,p_{n}$} be primes. Show that the expression $p_{1}p_{2}\ldots p_{n} + 1$ is divisible by none of these primes.
\begin{proof}
    Let $p\in \{p_{1}, p_{2}, \ldots, p_{n}\}$. By contradiction, assume
    \begin{align*}
        p \mid p_{1}p_{2}\ldots p_{n} + 1.
    \end{align*}
    This implies that $\exists \ s\in \mathbb{Z}$ s.t
    \begin{align*}
        & p_1p_2\ldots p_n + 1 = ps.
    \end{align*}
    But, since $p$ is prime, it's divisors can only be $1$ or $p$. Since every prime number is greater than 1, $p \nmid 1$. So, the only way for this expression to be true is if
    \begin{align*}
        p_1p_2\ldots p_n + 1 = p
    \end{align*}
    which we know cannot be true. $\therefore \ p \nmid p_1p_2\ldots p_n + 1$.
\end{proof}
\textbf{(b)} Use part (a) to prove that there are infinitely many primes.
\begin{proof}
    By contradiction, assume that there are finitely many primes. There must exist some maximum prime, call it $p_{f}$. By (a), 
    \begin{align*}
        &p_{1}, p_{2}, \ldots , p_{f} \nmid p_{1}p_{2}\ldots p_{f} + 1. 
    \end{align*}
    Let $m = p_{1}p_{2}\ldots p_{f} + 1 \in \mathbb{Z}$. By the Fundamental Theorem of Arithmetic, m is either a prime or a product of primes. But if m is a prime,
    \begin{align*}
        &p_{f} < p_{1}p_{2}\ldots p_{f} + 1 \\
        &p_{f} < m
    \end{align*}
    contradicts our assumption. So, $m$ must be a product of primes, which implies $\exists$ a prime $p_{i}$ that divides it. But that contradicts (a) and $\therefore m$ can not be a product of primes. Thus, this contradicts the Fundamental Theorem of Arithmetic, and shows that the assumption of $p_{f}$ being the largest prime must be false. $\therefore$ there must be infinitely many primes.
\end{proof}

\textbf{2.} For every positive integer $n$, prove that a set with exactly $n$ elements has exactly $2^{n}$ subsets (including the empty set and the entire set).
\begin{proof}
    We proceed by induction on $n$. \\
    \indent{\underline{Base Case}:} Let $n = 1$. Then, the set $\{1\}$ has the subsets
    \begin{align*}
        \underbrace{\emptyset, \{1\}}_{2^{1}}
    \end{align*}
    which amount to $2^{1}$ subsets. \\
    \indent{\underline{Induction Step}:} Let $n = k$ where k is some integer greater than $1$. Assume the set $S_{1} = \{1, 2, 3,\ldots, k\}$ has the subsets
    \begin{align*}
        \underbrace{\emptyset, \{1\}, \{2\}, \ldots, \{1, 2, 3,\ldots, k\}}_{2^{k}}
    \end{align*}
    which amount to $2^{k}$ subsets. We must show that for $n = k + 1$ the set $S_{2} = \{1, 2, 3,\ldots, k, k + 1\}$ has $2^{k+1}$ subsets. But, since $S_{1}$ is already contained in $S_{2}$, all of $S_{1}$'s subsets must also be a subset of $S_{2}$. Thus, the current number of subsets in $S_{2}$ is
    \begin{align*}
        2^{k} + \underbrace{\{1, 2, 3, \ldots, k, k+1\} + \ldots}_{\text{remaining amount of subsets}}.
    \end{align*}
    But, the only difference between the subsets of $S_{1}$ and the remaining subsets of $S_{2}$ is the element $k + 1$. So, the remaining subsets, excluding $k + 1$, must each be a subset of $S_{1}$ which we know has $2^{k}$ subsets. So the total subsets of $S_{2}$ must be
    \begin{align*}
        2^{k} + 2^{k} &= 2(2^{k}) \\
        &= 2^{k+1}.
    \end{align*}
    So, a set of $n$ elements has $2^{n}$ subsets for $n > 0$.
\end{proof}

\textbf{3.} The Fibonacci numbers are $1,1,2,3,5,8, 13, 21, 34, \ldots$. In general, the Fibonacci numbers are defined by $f_{1} = 1$, $f_{2} = 1$, and for $n \geq 3$, $f_{n} = f_{n-1} + f_{n-2}$. Prove that the $n$th Fibonacci number $f_{n}$ satisfies $f_{n} < 2^{n}$ for all $n \geq 1$.
\begin{proof}
    We proceed by induction on $n$. \\
    \indent{\underline{Base Case}:} Let $n=1$. Then,
    \begin{align*}
        f_{1} &= 1 \\
        \Rightarrow \ 1 &< 2^{1} = 2 \ \checkmark
    \end{align*}  \\
    \indent{\underline{Induction Step}:} Let $n=k$ where $k \in \mathbb{Z}$ and $k \geq 1$. Assume that $f_{k} < 2^{k}$. We must show that for $n = k + 1$, $f_{k+1} < 2^{k+1}$. By the definition of the Fibonacci numbers, we know that
    \begin{align*}
        f_{k + 1} = f_{k} + f_{k-1}.
    \end{align*}
    But, we know that $f_{k} < 2^{k}$. So,
    \begin{align*}
        f_{k + 1} < 2^{k} + f_{k-1}.
    \end{align*}
    But, $f_{k-1} < f_{k}$ which implies $f_{k-1} < 2^{k}$. Thus, 
    \begin{align*}
        f_{k + 1} &< 2^{k} + 2^{k} \\
        \Rightarrow \ f_{k + 1} &< 2(2^{k}) \\
        \Rightarrow \ f_{k + 1} &< 2^{k + 1}.
    \end{align*}
    $\therefore$ $f_{n} < 2^{n}$ holds $\forall$ $n \geq 1$
\end{proof}

\textbf{4.} Let $a,b \in \mathbb{R}$, define $a \sim b$ if $a-b$ is an integer. Show that $\sim$ is an equivelance relation on $\mathbb{R}$. Describe the equivalence classes produced by this relation. \\ \\
 To show that $\sim$ is an equivalence relation on $\mathbb{R}$, we must show that the properties of one hold. \\
 \indent{\underline{Reflexivity}:} Let $a \in \mathbb{R}$. We want to show that $a \sim a$. That is,
 \begin{align*}
     &a -a \in \mathbb{Z} \ ? \\
    \Rightarrow \ &0 \in \mathbb{Z} \ \checkmark
 \end{align*}
\indent{\underline{Symmetry}:} Suppose $a \sim b$. We want to show that $b \sim a$. That is,
\begin{align*}
    b - a \in \mathbb{Z} \ ?
\end{align*}
Since $a \sim b$, $a-b \in \mathbb{Z}$ $\therefore$ $b-a \in \mathbb{Z}$ but with the signs flipped. \\ \\
\indent{\underline{Transitivity}:} Suppose $a \sim b$ and $b \sim c$. We want to show that $a \sim c$. That is
\begin{align*}
    a - c \in \mathbb{Z} \ ?
\end{align*}
Since $a \sim b$ and $b \sim c$,
\begin{align*}
    &a-b, \ b-c \in \mathbb{Z}\\
    \Rightarrow \ &(a-b) + (b-c) \in \mathbb{Z} \\
    \Rightarrow \ &a-c \in \mathbb{Z} \ \checkmark
\end{align*}
Thus, $\sim$ is an equivalence relation on $\mathbb{R}$.
 Equivalence classes produced by $\sim$ are of the form 
\begin{align*}
    \underbrace{[a] = \{x.a \mid x,a \in \mathbb{R}, x \sim a\}}_{\text{the class of all real numbers with $a$ following the decimal point}}
\end{align*} 
\textbf{5.} Suppose $f : A \xrightarrow{} B$ and $g: B \xrightarrow{}C$ are two functions. \\
\indent{\textbf{(a)} If $f$ and $g$ are one-to-one, show that the function $gf$ (the composition $g \circ f$) is also one-to-one.
\begin{proof}
    Suppose $f$ and $g$ are one-to-one. Let $a_{1}, a_{2} \in A$ such that
    \begin{align*}
        gf(a_{1}) = gf(a_{2}).
    \end{align*}
    Since $g$ is one-to-one, this equality implies that
    \begin{align*}
       f(a_{1}) = f(a_2).
    \end{align*}
    But, $f$ is also one-to-one. So,
    \begin{align*}
        a_1 = a_2
    \end{align*}
    implies $g \circ f$ is one-to-one.
\end{proof}
\indent{\textbf{(b)}} If $f$ and $g$ are onto, show that the function $gf$ is also onto. 
\begin{proof}
    Suppose $f$ and $g$ are onto. Let $a_{1} \in A$ and $b_{1} \in B$ such that 
    \begin{align*}
        gf(a_1) = g(b_1).
    \end{align*}
    Since $f$ is onto, $\exists$ some $a_1 \in A$ s.t $f(a_1) = b_1$. But, $g$ is also onto, which implies that $\exists$ some $c_1 \in C$ s.t $g(b_1) = c_1$. Rewriting the composition,
    \begin{align*}
        gf(a_1) = c_1
    \end{align*}
    shows that $g \circ f$ is onto.
\end{proof}
\end{document}